\input{../tex-inputs/latex-leseansicht-vorspann}

\section[Carl von Torresani an Arthur Schnitzler, 8. 10. {[}1892{]}]{L04325 Carl von Torresani an Arthur Schnitzler, 8. 10. {[}1892{]}}
\nopagebreak\mylabel{L04325v}
\rehead{ }\normalsize\beginnumbering\briefempfaengerindex{Schnitzler, Arthur@\textsc{Schnitzler, Arthur}!zzzTorresani-Lanzenfeld, Carl von@\emph{von Carl von Torresani-Lanzenfeld}!1892-10-083@{8. 10. {[}1892{]}}|(be}
\toendnotes[C]{\smallbreak\pagebreak[2]}
\correspDesc{Versand  durch Carl von Torresani am 8. 10. [1892] in Wien
\newline{}Erhalt  durch Arthur Schnitzler im Zeitraum [8. 10. 1892
                  – 11. 10. 1892?] in Wien}\toendnotes[C]{\smallbreak}
\Standort{CUL, Schnitzler, B 106.}
\physDesc{Brief, 1 Blatt, 1 Seite, 301 Zeichen (Briefpapier mit Pflanzenmotiv (Heidelbeerzweig))
\newline{}Handschrift: schwarze Tinte, deutsche Kurrent
\newline{}Schnitzler: mit Bleistift Jahresangabe bei der Datierung ergänzt: »9\textcolor{gray}{2}« }\toendnotes[C]{\smallbreak}
\pstart
           \centering{}{\pb}Wien\oindex{Wien@\textbf{Wien}, \emph{Verwaltungsgebiet}|pw}{ }8/10\pend
           
\pstart{}Verehrteſter Herr Doktor!\pend\vspace{0.5em}
\pstart
           Ich kann zu meinem Leidweſen morgen Sonntag{ }\label{K_L04325-1v}\edtext{Ihrer fdl Aufforderung\eventindex{Wollzeile 15 (»Berthahof«)@\textbf{Wollzeile 15 (»Berthahof«)}!Private Lesung von Gestern, 9.10.1892@Private Lesung von Gestern, 9.10.1892|pwv}}{\lemma{\textnormal{\emph{Ihrer fdl Aufforderung}}}\Cendnote{\textnormal{Schnitzlers Einladung ist nicht überliefert,
                  aber sein Brief an Salten\pwindex{Salten, Felix 6.\,9.\,1869 Budapest – 8.\,10.\,1945 Zürich@\textsc{Salten, Felix} (6.\,9.\,1869 Budapest – 8.\,10.\,1945 Zürich), \emph{Schriftsteller, Journalist, Chefredakteur}|pwk}, in dem er zu einem Besuch der \emph{Ausstellung für Musik und Theaterwesen}\orgindex{Internationale Ausstellung für Musik und Theaterwesen@Internationale Ausstellung für Musik und Theaterwesen|pwk} gemeinsam mit Beer-Hofmann\pwindex{Beer-Hofmann, Richard 11.\,7.\,1866 Wien – 26.\,9.\,1945 New York City@\textsc{Beer-Hofmann, Richard} (11.\,7.\,1866 Wien – 26.\,9.\,1945 New York City), \emph{Schriftsteller}|pwk} anregte
                  und berichtete, er habe auch Torresani\pwindex{Torresani-Lanzenfeld, Carl von 19.\,4.\,1846 Mailand – 16.\,4.\,1907 Nago-Torbole@\textsc{Torresani-Lanzenfeld, Carl von} (19.\,4.\,1846 Mailand – 16.\,4.\,1907 Nago-Torbole), \emph{Schriftsteller, Offizier}|pwk} informiert, siehe XXXX Auszeichnungsfehler: Dokument L02957 nicht gefunden. Auch Salten\pwindex{Salten, Felix 6.\,9.\,1869 Budapest – 8.\,10.\,1945 Zürich@\textsc{Salten, Felix} (6.\,9.\,1869 Budapest – 8.\,10.\,1945 Zürich), \emph{Schriftsteller, Journalist, Chefredakteur}|pwk} sagte ab (XXXX Auszeichnungsfehler: Dokument L03115 nicht gefunden). Es fand stattdessen am Sonntagnachmittag eine Deklamation\eventindex{Wollzeile 15 (»Berthahof«)@\textbf{Wollzeile 15 (»Berthahof«)}!Private Lesung von Gestern, 9.10.1892@Private Lesung von Gestern, 9.10.1892|pwk} von Hofmannsthals\pwindex{Hofmannsthal, Hugo von 1.\,2.\,1874 Wien – 15.\,7.\,1929 Rodaun@\textsc{Hofmannsthal, Hugo von} (1.\,2.\,1874 Wien – 15.\,7.\,1929 Rodaun), \emph{Schriftsteller}|pwk} Einakter \emph{Gestern}\pwindex{Hofmannsthal, Hugo von 1.\,2.\,1874 Wien – 15.\,7.\,1929 Rodaun@\textsc{Hofmannsthal, Hugo von} (1.\,2.\,1874 Wien – 15.\,7.\,1929 Rodaun), \emph{Schriftsteller}!Gestern. Dramatische Studie in einem Akt in Versen@\strich\emph{Gestern. Dramatische Studie in einem Akt in Versen}|pwk} bei Beer-Hofmann\pwindex{Beer-Hofmann, Richard 11.\,7.\,1866 Wien – 26.\,9.\,1945 New York City@\textsc{Beer-Hofmann, Richard} (11.\,7.\,1866 Wien – 26.\,9.\,1945 New York City), \emph{Schriftsteller}|pwk}
                  statt, wie aus dem \emph{Tagebuch}\pwindex{Schnitzler, Arthur 15. 5. 1862 Wien – 21. 10. 1931 ebd.@\textsc{Schnitzler, Arthur} (15. 5. 1862 Wien – 21. 10. 1931 ebd.), \emph{Schriftsteller, Mediziner}!Tagebuch@\strich\emph{Tagebuch}|pwk} hervorgeht: »Richard\pwindex{Beer-Hofmann, Richard 11.\,7.\,1866 Wien – 26.\,9.\,1945 New York City@\textsc{Beer-Hofmann, Richard} (11.\,7.\,1866 Wien – 26.\,9.\,1945 New York City), \emph{Schriftsteller}|pw} im Tricot, ich im
                     Schlafrock«, siehe A. S.: \emph{Tagebuch}, 9. 10. 1892.}}}\label{K_L04325-1} nicht nachko{\geminationm}en, da ich
               bei meiner Schwiegermutter\pwindex{Pabst, Amalia @\textsc{Pabst, Amalia}|pwv}
               in Baden\oindex{Baden bei Wien@\textbf{Baden bei Wien}, \emph{Hauptstadt}|pw} dinire. Doch werde ich gegen halb eilf
               Nachts \label{K_L04325-2v}\edtext{zum \textsc{Kremser\oindex{Wien@\textbf{Wien}!I., Innere Stadt@\textbf{I., Innere Stadt}!Café Kremser@\textbf{Café Kremser}, \emph{Kaffeehaus}|pw}}}{\lemma{\textnormal{\emph{zum Kremser}}}\Cendnote{\textnormal{Schnitzlers{ }\emph{Tagebuch}\pwindex{Schnitzler, Arthur 15. 5. 1862 Wien – 21. 10. 1931 ebd.@\textsc{Schnitzler, Arthur} (15. 5. 1862 Wien – 21. 10. 1931 ebd.), \emph{Schriftsteller, Mediziner}!Tagebuch@\strich\emph{Tagebuch}|pwk} bestätigt ein Treffen im Kaffeehaus\oindex{Wien@\textbf{Wien}!I., Innere Stadt@\textbf{I., Innere Stadt}!Café Kremser@\textbf{Café Kremser}, \emph{Kaffeehaus}|pwk}. Neben Torresani\pwindex{Torresani-Lanzenfeld, Carl von 19.\,4.\,1846 Mailand – 16.\,4.\,1907 Nago-Torbole@\textsc{Torresani-Lanzenfeld, Carl von} (19.\,4.\,1846 Mailand – 16.\,4.\,1907 Nago-Torbole), \emph{Schriftsteller, Offizier}|pwk} war auch Hermann Bahr\pwindex{Bahr, Hermann 19.\,7.\,1863 Linz – 15.\,1.\,1934 München@\textsc{Bahr, Hermann} (19.\,7.\,1863 Linz – 15.\,1.\,1934 München), \emph{Schriftsteller, Kritiker}|pwk} anwesend, siehe A. S.: \emph{Tagebuch}, 9. 10. 1892.}}}\label{K_L04325-2} kommen u hoffe die Herren\pwindex{Beer-Hofmann, Richard 11.\,7.\,1866 Wien – 26.\,9.\,1945 New York City@\textsc{Beer-Hofmann, Richard} (11.\,7.\,1866 Wien – 26.\,9.\,1945 New York City), \emph{Schriftsteller}|pwv}\pwindex{Hofmannsthal, Hugo von 1.\,2.\,1874 Wien – 15.\,7.\,1929 Rodaun@\textsc{Hofmannsthal, Hugo von} (1.\,2.\,1874 Wien – 15.\,7.\,1929 Rodaun), \emph{Schriftsteller}|pwv} dort zu
               treffen.\pend
           
\pstart
           In herzl. Ergebenheit Ihr Sie hochachtender{\\[\baselineskip]}\spacefill\mbox{Torresani}\pend
           \leftskip=0em{}\selectlanguage{ngerman}\endnumbering\briefempfaengerindex{Schnitzler, Arthur@\textsc{Schnitzler, Arthur}!zzzTorresani-Lanzenfeld, Carl von@\emph{von Carl von Torresani-Lanzenfeld}!1892-10-083@{8. 10. {[}1892{]}}|)be}\mylabel{L04325h}
\begin{anhang}
\end{anhang}\newcommand{\dateiname}{L04325}\newcommand{\titel}{Carl von Torresani an Arthur Schnitzler, 8. 10. [1892]}\newcommand{\editorInnen}{Herausgegeben von Jahnke, SelmaMüller, Martin Anton}\input{../tex-inputs/latex-leseansicht-abspann}
